% Appendix C -- The simplified model of Section 3.3: proofs (draft prose)
%   source: model-codes/codes/epc-lifecycle-deep-learning/notes/
%           stylized-mechanism-2026-10-16/phase-regions.tex, Appendix A
%   timing  estimation_epc/model_funcs_estimation_epc.py: bequest_value
%           (estate at the end-of-year price), reward (util + (1-surv) g),
%           renovate in place pays the inefficient bill that year,
%           renovate-at-buy and buy-efficient pay the efficient bill at once
\section{The Simplified Model: Proofs}\label{app:proofs}

\paragraph{Setup.}
We use the simplified model of Section~\ref{sec:model-forces}. We suppress
the parts of the state that do not change, and write $V_{j+1}(a,t)$ for the
value next year, before the taste shocks are drawn, of an owner of a home
of type $t$ that leaves this year with financial wealth $a$, saving less
mortgage. Next year the owner draws one shock for keeping its home and one
for selling it, distributed as in Section~\ref{sec:model-problem}. Let
$Q^k_{j+1}(a,t)$ be the value of keeping the home: that of Stay for an
efficient home, and the higher of $Q^{\mathrm{stay},N}_{j+1}$ and
$Q^{\mathrm{stay},R}_{j+1}$ for an inefficient one. Let $Q^m_{j+1}(a,t)$ be the value of selling,
the highest value among the alternatives that sell the home. Then
\begin{equation}
  V_{j+1}=\sigma_\varepsilon\log\Bigl[\exp\bigl(Q^k_{j+1}/\sigma_\varepsilon\bigr)
  +\exp\bigl(Q^m_{j+1}/\sigma_\varepsilon\bigr)\Bigr],
  \label{eq:app-logsum}
\end{equation}
and the owner keeps its home with probability
$\partial V_{j+1}/\partial Q^k_{j+1}$. $V'_{j+1}$ is the marginal value of
cash on hand next year, so that $\partial V_{j+1}/\partial a=(1+r)V'_{j+1}$,
where $r$ is the rate on the asset the household adjusts at the margin; we
write $V'^{\,k}_{j+1}$ and $V'^{\,m}_{j+1}$ for its values in the two
branches. When the household can adjust
that asset without hitting a borrowing limit, its Euler equation is
\begin{equation}
  u_c=\beta\phi_j(1+r)V'_{j+1}+(1-\phi_j)\,g',
  \label{eq:app-euler}
\end{equation}
since by equation~\eqref{eq:end-value} the estate is not discounted. With
$\tilde\phi_j$ from equation~\eqref{eq:phi-tilde}, this gives
$\tilde\phi_j=\beta\phi_j(1+r)V'_{j+1}/u_c$,
$1-\tilde\phi_j=(1-\phi_j)g'/u_c$, and
$\tilde\phi_j\Lambda=\beta\phi_j\varrho V'_{j+1}/u_c$. We write $R_s\equiv(1-\delta_h-\zeta_s)\Delta p$ for what the efficient
state adds to the proceeds of a sale, and $q\equiv\tilde\phi\Lambda$.

\paragraph{Outline.}
The proof of Proposition~\ref{prop:inplace} compares renovating this year
with renovating next year. Lemma~\ref{lem:decomp} writes the gain from
renovating now as this year's comfort, less the cost, plus the value of the
efficient state next year. Lemma~\ref{lem:match} bounds that value. An owner
of the inefficient home can match an efficient owner next year: if the
efficient owner keeps its home, by renovating, which costs $\overline C$
more; if it sells, by selling too, with $R_s$ less. Whether an owner sells
can itself depend on whether its home is efficient. But the owners whose
choice turns on efficiency are close to indifferent between keeping and
selling, so to first order the state is worth the same to them either way.

The proof of Proposition~\ref{prop:buy} compares the buyer with an owner of
the inefficient home who lives as the buyer does, but without the state. It
needs no proviso about renovating later. The state is cheaper to buy than
to make, $\overline C_b<C(h')+\Delta e'-\delta_h\Delta p'$, so near the
buyer's threshold an owner of the inefficient home would not renovate it.

\paragraph{Two routes, one horizon.}
Let $\mathcal C$ be the extra cost in the year the household gets the
efficient state, over the extra cost of holding it for a year:
$\mathcal C=\overline C$ in place and $\mathcal C=\overline C_b$ at a
purchase. Rearranged, conditions \eqref{eq:rule-inplace}
and~\eqref{eq:rule-buy} both say that the household gets the efficient
state if and only if
\begin{equation}
  \theta c+(\Delta e-\delta_h\Delta p)+(1-\tilde\phi)\,\Delta p
  +\tilde\phi\,(1-\mu)\,\Lambda\,(1-\delta_h-\zeta_s)\Delta p
  \;\ge\;\frac{\mathcal C}{\widetilde H},
  \qquad
  \widetilde H\equiv\frac{1}{1-\tilde\phi\Lambda\mu},
  \label{eq:common}
\end{equation}
with $\theta$, $\Delta e$, and $\Delta p$ those of the home concerned. The
left side is what the efficient state yields in a year: the comfort, the
bill saving net of the maintenance on the premium, the premium that the
heirs receive if the owner dies, and the premium net of costs that the
owner receives if it moves. The right side spreads the cost $\mathcal C$
over the horizon $\widetilde H$. At constant parameters
$H_j=1/(1-\beta\phi\varrho\mu)$ in equation~\eqref{eq:payoff-horizon},
while $\tilde\phi\Lambda=\beta\phi\varrho V'_{j+1}/u_c$ by the Euler
equation. So $\widetilde H$ is $H_j$ with next year weighted by its
marginal utility relative to this year. The two routes face the same
horizon, but spread different costs over it.

\begin{lemma}[Decomposition]\label{lem:decomp}
Take an owner of an inefficient home of size $h$ at age $j$ that keeps its
home this year. Let $c$ and $a'$ be its consumption and end-of-year
wealth if it does not renovate, $C\equiv C(h)$, and suppose that it can reduce $a'$ by
$C$ at the rate $r$. Then $G_j\equiv Q_j^{\mathrm{stay},R}-Q_j^{\mathrm{stay},N}$
satisfies, to first order,
\begin{equation}
  \frac{G_j}{u_c}\approx\theta c-C+\tilde\phi_j\Lambda D_{j+1}
  +(1-\tilde\phi_j)\Delta p,
  \qquad
  D_{j+1}\equiv\frac{V_{j+1}(a'-C,R)-V_{j+1}(a'-C,N)}{V'_{j+1}},
  \label{eq:app-decomp}
\end{equation}
where $D_{j+1}$ is the money value of being efficient next year.
\end{lemma}

\begin{proof}
If it does not renovate, $Q_j^{\mathrm{stay},N}=u_j(c,s_N)+(1-\phi_j)g(a'+p^Nh)
+\beta\phi_jV_{j+1}(a',N)$. Renovating, consuming $c$, and leaving the year
with $a'-C$ is feasible. In the year of the work the renovator receives
$s_R$ and pays the inefficient bill and maintenance
(Appendix~\ref{app:model-budgets}). If it dies, the estate values the home
at $p^Rh$. If it survives, the home is efficient next year with probability
$\varrho$. So
\begin{multline*}
  G_j\ge\bigl[u_j(c,s_R)-u_j(c,s_N)\bigr]
  -\bigl[(1-\phi_j)\bigl(g(a'+p^Nh)-g(a'-C+p^Nh)\bigr)
  +\beta\phi_j\bigl(V_{j+1}(a',N)-V_{j+1}(a'-C,N)\bigr)\bigr]\\
  +\beta\phi_j\varrho\bigl[V_{j+1}(a'-C,R)-V_{j+1}(a'-C,N)\bigr]
  +(1-\phi_j)\bigl[g(a'-C+p^Rh)-g(a'-C+p^Nh)\bigr].
\end{multline*}
To first order the four terms are $u_c\theta c$, $-Cu_c$ by the Euler
equation~\eqref{eq:app-euler}, $\beta\phi_j\varrho V'_{j+1}D_{j+1}
=u_c\tilde\phi_j\Lambda D_{j+1}$, and $(1-\phi_j)g'\Delta p
=u_c(1-\tilde\phi_j)\Delta p$. The renovator's best allocation differs from
$(c,a'-C)$ by a first-order amount, so by the envelope theorem the
inequality holds with equality to first order.
\end{proof}

\begin{lemma}[Matching]\label{lem:match}
Let
\begin{equation*}
  D^k\equiv\frac{Q^k_{j+1}(a,R)-Q^k_{j+1}(a,N)}{V'^{\,k}_{j+1}},
  \qquad
  D^m\equiv\frac{Q^m_{j+1}(a,R)-Q^m_{j+1}(a,N)}{V'^{\,m}_{j+1}}
\end{equation*}
be the money values of the efficient state next year to an owner that
keeps its home and to one that sells it. To first order:
\begin{enumerate}[label=(\roman*)]
\item $D_{j+1}=\mu_{j+1}D^k+(1-\mu_{j+1})D^m$, where $\mu_{j+1}$ is the
  probability of keeping the home, averaged between the inefficient and
  the efficient state and weighted by marginal utility;
\item $D^k\le\overline C$, with equality if an owner of the inefficient
  home that keeps it would renovate it;
\item $D^m=R_s$;
\item $D_{j+1}\le U_{j+1}$, the money value of the state along the
  efficient owner's path: the comfort and the bill saving net of
  maintenance in each year it holds the state, the premium its heirs
  receive if it dies, and $R_s$ if it sells. Equality holds if an owner of
  the inefficient home would never renovate it.
\end{enumerate}
\end{lemma}

\begin{proof}
(i) By \eqref{eq:app-logsum}, the derivatives of $V_{j+1}$ with respect to
$Q^k_{j+1}$ and $Q^m_{j+1}$ are the probabilities of keeping and of selling
the home. Integrating along the line from the values of the inefficient
state to those of the efficient one gives, exactly,
\begin{equation*}
  V_{j+1}(a,R)-V_{j+1}(a,N)=\bar\mu\bigl[Q^k_{j+1}(a,R)-Q^k_{j+1}(a,N)\bigr]
  +(1-\bar\mu)\bigl[Q^m_{j+1}(a,R)-Q^m_{j+1}(a,N)\bigr],
\end{equation*}
where $\bar\mu$ is the probability of keeping the home averaged along the
line. Dividing by $V'_{j+1}$ gives (i), with
$\mu_{j+1}=\bar\mu V'^{\,k}_{j+1}/\bigl[\bar\mu V'^{\,k}_{j+1}+(1-\bar\mu)V'^{\,m}_{j+1}\bigr]$;
the denominator differs from $V'_{j+1}$ by a first-order amount, which
multiplies a first-order term. The probabilities of keeping the home in the
two states also differ by a first-order amount, so to first order it does
not matter at which state they are taken. (ii) An owner of the
inefficient home that keeps it can renovate and live as an efficient owner
that stays. It gets the same services $s_R$ and the same home next year,
and its borrowing limit is no tighter, since it may borrow for the work.
It pays $C+(\delta_hp^N+e^N)h$ instead of $(\delta_hp^R+e^R)h$, that is,
$\overline C$ more. So $D^k\le\overline C$, with equality when this is its
best choice. (iii) An efficient home adds $R_s$ to the proceeds of a sale
(Appendix~\ref{app:model-budgets}), and after the sale the two owners face
the same choices. (iv) An owner of the inefficient home can follow the
efficient owner's path: given the same taste shocks, it keeps or sells its
home when the efficient owner does, with the same consumption, but without
the state. Its taste shocks are then the same, so to first order this
costs it $U_{j+1}$, computed with the efficient owner's probabilities of
keeping its home. So $D_{j+1}\le U_{j+1}$, with equality when never
renovating is its best choice; its own choices between keeping and selling
then differ from the efficient owner's only at the margin, where they cost
nothing to first order.
\end{proof}

\begin{proof}[Proof of Proposition~\ref{prop:inplace}]
Let $\ell_{j+1}\equiv\overline C-D_{j+1}$. Substituting
$D_{j+1}=\overline C-\ell_{j+1}$ in \eqref{eq:app-decomp},
\begin{equation}
  \frac{G_j}{u_c}\approx\theta c-\bigl(C-\Lambda\overline C\bigr)
  -(1-\tilde\phi_j)\bigl(\Lambda\overline C-\Delta p\bigr)
  -\tilde\phi_j\Lambda\,\ell_{j+1}.
  \label{eq:app-gap}
\end{equation}
By Lemma~\ref{lem:match}(i)--(iii),
$\ell_{j+1}=\mu_{j+1}(\overline C-D^k)+(1-\mu_{j+1})\ell\ge(1-\mu_{j+1})\ell$,
so $G_j\ge0$ implies \eqref{eq:rule-inplace}. Under the proviso,
$D^k=\overline C$ and $\ell_{j+1}=(1-\mu_{j+1})\ell$, so \eqref{eq:app-gap}
has the sign of \eqref{eq:rule-inplace}, which is then also sufficient.
\end{proof}

\begin{remark}[Constant parameters]\label{rem:constant}
Suppose that $\tilde\phi$, $\mu$, $\Lambda$, and consumption are constant
over the following years, and that no borrowing limit binds. Then
\eqref{eq:rule-inplace} is necessary and sufficient without the proviso.
Let $Z\equiv\theta c+(\Delta e-\delta_h\Delta p)+(1-\tilde\phi)\Delta p$,
the yield of the state in a year in which the owner keeps it. Since
$\theta c-C+(1-\tilde\phi)\Delta p=Z-\overline C$, \eqref{eq:app-decomp}
reads $G_j/u_c\approx Z-\overline C+qD_{j+1}$. Next year, an owner of the
inefficient home can renovate the first time it keeps its home, which values the
state at $\overline D\equiv\mu\overline C+(1-\mu)R_s$; never renovate,
which values it at $U$ with $U=\mu(Z+qU)+(1-\mu)R_s$; or renovate after
$k$ years, which values it at $U+(\mu q)^k(\overline D-U)$. It takes the
best of these, so $D_{j+1}=\min\{U,\overline D\}$. Let
$E\equiv Z-(1-\mu q)\overline C+q(1-\mu)R_s$; $E\ge0$ is
\eqref{eq:common} with $\mathcal C=\overline C$, that is,
\eqref{eq:rule-inplace} rearranged. If $D_{j+1}=\overline D$, then
$G_j/u_c\approx E$. If $D_{j+1}=U$, then
$(1-\mu q)\,G_j/u_c\approx(1-\mu q)(Z-\overline C)+q\bigl[\mu Z+(1-\mu)R_s\bigr]=E$.
Either way $G_j$ has the sign of $E$.
\end{remark}

\begin{proof}[Proof of Proposition~\ref{prop:buy}]
Write $R_s'\equiv(1-\delta_h-\zeta_s)\Delta p'$,
$Z'\equiv\theta'c+(\Delta e'-\delta_h\Delta p')+(1-\tilde\phi)\Delta p'$,
and $\overline C'\equiv C(h')+\Delta e'-\delta_h\Delta p'$ for the in-place
cost of $h'$. Compare the two purchases at the consumption $c$ and the size
$h'$ of the buyer of the inefficient home. Getting the state costs $C_b$
more this year and saves the bill $\Delta e'$, so the buyer of the state
leaves the year with $C_b-\Delta e'$ less. If it dies, the estate values
the home at $p^Rh'$. As in Lemma~\ref{lem:decomp}, to first order the gain
$G_b\equiv Q_j^{\mathrm{buy},R}-Q_j^{\mathrm{buy},N}$ satisfies
\begin{equation*}
  \frac{G_b}{u_c}\approx\theta'c+\Delta e'-C_b+qD'+(1-\tilde\phi)\Delta p'
  =Z'-\overline C_b+qD',
\end{equation*}
where $D'$ is the money value of the state next year. As in
Remark~\ref{rem:constant}, an owner of the inefficient home $h'$ can
renovate it the first time it keeps it, which values the state at
$\overline D'\equiv\mu\overline C'+(1-\mu)R_s'$, or never renovate it, which
values it at $U'=[\mu Z'+(1-\mu)R_s']/(1-\mu q)$, or wait; so
$D'=\min\{U',\overline D'\}$. Both $Z'$ and $U'$ rise with $c$, so $G_b$ is
continuous and strictly increasing in $c$. Where $D'=U'$,
\begin{equation*}
  (1-\mu q)\,\frac{G_b}{u_c}\approx
  E_b\equiv Z'-(1-\mu q)\overline C_b+q(1-\mu)R_s',
\end{equation*}
and $E_b\ge0$ is \eqref{eq:common} with $\mathcal C=\overline C_b$, that
is, \eqref{eq:rule-buy} rearranged. At the consumption where $E_b=0$,
$U'=\mu\overline C_b+(1-\mu)R_s'\le\overline D'$, since
$\overline C_b\le C(h')-\delta_h\Delta p'<\overline C'$. So $D'=U'$ and
$G_b=0$ there. As $G_b$ and $E_b$ both increase in $c$, $G_b\ge0$ if and
only if $E_b\ge0$.
\end{proof}

\begin{lemma}[The payoff horizon]\label{lem:horizon}
Let a household buy a home of size $h'$ at age $j$, and let
$\tilde\phi_{j+k}$, $\Lambda_{j+k}$, $\mu_{j+k+1}$, and $c_{j+k}$ vary with
$k$. Suppose that an owner of the inefficient home $h'$ would never
renovate it, and that no borrowing limit binds. Let $P_0=1$ and
$P_k=\prod_{i<k}\tilde\phi_{j+i}\Lambda_{j+i}\mu_{j+i+1}$, the discounted
probability of still holding the state $k$ years after the purchase. Then,
to first order, the household buys the home efficient if and only if
\begin{equation*}
  \sum_{k\ge0}P_k\Bigl[\theta'c_{j+k}+(\Delta e'-\delta_h\Delta p')
  +(1-\tilde\phi_{j+k})\Delta p'
  +\tilde\phi_{j+k}\Lambda_{j+k}(1-\mu_{j+k+1})R_s'\Bigr]
  \;\ge\;\overline C_b.
\end{equation*}
The weights sum to $\widetilde H_j\equiv\sum_kP_k$. Since
$\tilde\phi_{j+i}\Lambda_{j+i}=\beta\phi_{j+i}\varrho V'_{j+i+1}/u_{c,j+i}$
and $V'_{j+i+1}\approx u_{c,j+i+1}$,
\begin{equation*}
  P_k\approx\beta^k\Bigl(\prod_{i<k}\phi_{j+i}\,\varrho\,\mu_{j+i+1}\Bigr)
  \frac{u_{c,j+k}}{u_{c,j}},
\end{equation*}
so $\widetilde H_j$ is the payoff horizon $H_j$ of
equation~\eqref{eq:payoff-horizon}, with each year weighted by its marginal
utility relative to the year of purchase. At constant parameters the
condition is \eqref{eq:common} with $\mathcal C=\overline C_b$.
\end{lemma}

\begin{proof}
Write $Z_{j+k}$ for the first three terms in the bracket and
$q_{j+k}\equiv\tilde\phi_{j+k}\Lambda_{j+k}$. As in the proof of
Proposition~\ref{prop:buy}, $G_b/u_c\approx Z_j-\overline C_b+q_jD'_{j+1}$,
and under the assumptions $D'_{j+1}=U'_{j+1}$ by
Lemma~\ref{lem:match}(iv), where
$U'_{j+k}=\mu_{j+k}\bigl(Z_{j+k}+q_{j+k}U'_{j+k+1}\bigr)+(1-\mu_{j+k})R_s'$.
Substituting forward, and since $P_k\to0$,
$q_jU'_{j+1}=\sum_{k\ge1}P_kZ_{j+k}+\sum_{k\ge0}P_kq_{j+k}(1-\mu_{j+k+1})R_s'$.
Adding $Z_j-\overline C_b$ gives the condition.
\end{proof}

\begin{proof}[Proof of Proposition~\ref{prop:window}]
Under the assumptions of Proposition~\ref{prop:buy}, the in-place condition
\eqref{eq:rule-inplace} is necessary and sufficient by
Remark~\ref{rem:constant}. So the thresholds $c_j^*$ and $c_b^*$ are the
consumption at which \eqref{eq:common} holds with equality, with
$\mathcal C=\overline C$ and $\mathcal C=\overline C_b$, and
$\theta c^*=(1-\tilde\phi\Lambda\mu)\,\mathcal C-(\Delta e-\delta_h\Delta p)
-(1-\tilde\phi)\Delta p-\tilde\phi(1-\mu)\Lambda R_s$.
A higher probability of moving lowers $\mu$, and a higher probability of
dying lowers $\tilde\phi$.

\emph{Moving.} The derivative of $c^*$ in $1-\mu$ is
$\tilde\phi\Lambda(\mathcal C-R_s)/\theta$, and $\mathcal C-R_s$ is $\ell$
in place and $\ell_b$ at a purchase. Since $C_b\le(1+\zeta_b)\Delta p'$,
$\ell_b\le(\zeta_b+\zeta_s)\Delta p'$.

\emph{Death.} The derivative of $c^*$ in $1-\tilde\phi$ is
$[\Lambda(\mu\mathcal C+(1-\mu)R_s)-\Delta p]/\theta$. Here
$\Lambda(\mu\mathcal C+(1-\mu)R_s)$ is what the state is worth to the owner
alive next year, by Lemma~\ref{lem:match} in place and by the proof of
Proposition~\ref{prop:buy} at a purchase, and $\Delta p$ is what the heirs
receive. In place the derivative is at least
$(\Lambda\mu\overline C-\Delta p)/\theta$, which is positive under the
condition of part~(a). At a purchase,
$\mu\overline C_b+(1-\mu)R_s'\le(1+\zeta_b-\delta_h)\Delta p'$, since
$\overline C_b\le(1+\zeta_b-\delta_h)\Delta p'$ and $R_s'<\Delta p'$. So
the derivative is at most $[\Lambda(1+\zeta_b-\delta_h)-1]\Delta p'/\theta'$,
which is negative under the condition of part~(b). As $\Lambda\le\varrho$,
that condition holds when $\varrho<1/(1+\zeta_b-\delta_h)=0.981$.

\emph{The gap.} With $h'=h$,
$\theta(c^*_j-c^*_b)=(\overline C-\overline C_b)/\widetilde H$ and
$\overline C-\overline C_b=C(h)-C_b+\Delta e>0$, since $C_b\le C(h)$. And
$1/\widetilde H=1-\tilde\phi\Lambda\mu$ falls in $\tilde\phi$ and $\mu$. A
household with $c_b^*\le c<c_j^*$ would not renovate in place by
Proposition~\ref{prop:inplace}, and would take an efficient home at a
purchase by Proposition~\ref{prop:buy}.
\end{proof}
